\newpage

\begin{abstract}
Работа посвящена приложениям алгоритма Frank--Wolfe к задачам машинного обучения и смежным задачам крупномасштабной выпуклой оптимизации.
Цель работы состоит в развитии проекционно-свободных модификаций Frank--Wolfe, уменьшающих вычислительную стоимость итерации без потери базовых гарантий сходимости.
Для достижения этой цели рассматриваются методы с памятью о предыдущих направлениях, перенос CFW и NFW на гладкие задачи машинного обучения, стохастический блочно-координатный SOFW для транспортных сетей и децентрализованный блочно-координатный DBCFW.
Получены оценки сходимости для CFW, SOFW и DBCFW, описаны специальные случаи этих методов и проведены эксперименты для логистической регрессии и задач распределения транспортных потоков.
Результаты работы могут использоваться при выборе модификации Frank--Wolfe под конкретное вычислительное ограничение: стоимость линейного оракула, выбор направления или распределенный характер данных.

\smallskip
\textit{Ключевые слова}: Frank--Wolfe; условный градиент; машинное обучение; выпуклая оптимизация; моментум по направлениям; сопряженные направления; SOFW; стохастический LMO; блочно-координатная оптимизация; распределенная оптимизация.
\end{abstract}
